\documentclass[11pt]{article}
\usepackage{fontspec}
\setmainfont{Times New Roman}
\setsansfont{Arial}
\setmonofont{Consolas}
\usepackage{amsmath,amssymb}
\usepackage{geometry}
\usepackage{microtype}
\geometry{a4paper,margin=3cm}
\setlength{\parindent}{1.25em}
\setlength{\parskip}{0pt}
\emergencystretch=10em
\allowdisplaybreaks
\providecommand{\Kfield}{\mathfrak{K}}
\providecommand{\Ssys}{\mathfrak{S}}
\providecommand{\Mfield}{\mathfrak{M}}
\providecommand{\tq}{\tau}
\providecommand{\ds}{\displaystyle}

\begin{document}

\begin{center}
{\Large\bfseries 6. Поля и системы рациональных функций}\par
\vspace{1.2em}
Math. Ann. 76 (1915), S. 161--196
\end{center}

\vspace{2em}

Настоящая работа рассматривает вопросы о базах для произвольных систем рациональных и целых рациональных функций. Используемые здесь методы теории полей показывают, что существенным является решение этих вопросов для полей из рациональных функций --- полей рациональных функций\footnote{Ср. предварительное сообщение по случаю Венского съезда естествоиспытателей: Rationale Funktionenkörper --- Jahresber. d. D. Math.-Ver. 22 (1913) ---, где дан обзор постановок вопросов и результатов, относящихся к функциональным полям.} ---, тогда как обобщение результатов на произвольные системы получается как следствие.

Из вопросов о базах для общих систем до сих пор было известно лишь существование модульной базы, обеспеченное теоремой Гильберта (Math. Ann. 36). Ниже вопрос о рациональной представимости полностью решается через существование рациональной базы для всякой произвольной системы (\S\ 7); рациональная база полей получается уже в \S\ 4. Это существование рациональной базы позволяет всюду исходить из абстрактно определенного поля или системы и тем самым избегать трудностей, которые вызваны только специальным выбором рациональной базы, а не самой системой, например появлением специальных знаменателей или фундаментальных точек базисных функций.

Для полей будет также рассмотрен вопрос о минимальной базе, то есть о рациональной базе из алгебраически независимых функций (\S\ 6). Методы теории полей далее приводят к выделенной рациональной базе, инволюционной базе\footnote{Название выбрано потому, что рациональное функциональное поле в геометрическом толковании задает самую общую инволюцию в линейном пространстве.} (\S\ 5), которая становится особенно существенной для интегральных областей из полиномов (\S\ 8). Здесь она дает представление с фиксированным знаменателем (степенью функции, определенной интегральной областью), как это известно, например, в специальном случае типического представления инвариантов.

Из рациональной представимости теперь выводятся как можно более далеко идущие следствия для целой рациональной представимости, то есть для вопроса о конечности в более узком смысле.

Все известные до сих пор теоремы о конечности опираются на тот факт, что теорема Гильберта о модульной базе гарантирует конечность для всех таких систем, в которых представление
\[
  F=A_1f_1+A_2f_2+\cdots+A_kf_k
\]
влечет второе представление
\[
  F=B_1f_1+B_2f_2+\cdots+B_kf_k,
\]
где $B_i$ принадлежат системе и имеют меньшую степень, чем $F$.

Напротив, теорема о рациональной базе и методы теории полей позволяют заключать о конечности при предположениях совсем другого рода.

Так, как частичный ответ на одну из проблем Гильберта (Mathematische Probleme 14), получается класс относительно целых функций (относительно целые области первого типа), которые являются конечными интегральными областями и чья интегральная база задается упомянутой выше инволюционной базой (\S\ 10). По аналогии с гильбертовым доказательством теоремы Кронекера о фундаментальной системе алгебраически целых величин (Volle Invariantensysteme, \S\ 2; Math. Ann. 42) далее можно показать, что все регулярные системы из полиномов --- то есть все системы, которые можно отобразить в системы без фундаментальных точек, --- обладают интегральной базой (\S\ 12), тогда как нерегулярные системы без интегральной базы легко указать. В качестве простого примера этого факта упомянем теорему: ``Во всякой произвольной системе из бесконечно многих полиномов от одной неопределенной существует конечное число этих полиномов такое, что всякий полином системы можно представить как целую рациональную комбинацию этого конечного числа.'' Дальнейшие примеры находятся в \S\ 13.

Наконец, последние параграфы (\S\S\ 14 и 15) дают усиление теорем о базах в отношении целоалгебраичности.

Существенным вспомогательным средством исследования служит принцип переноса из \S\ 3, согласно которому достаточно рассматривать все поставленные здесь вопросы о базах для таких систем, в которых число неопределенных совпадает с числом алгебраически независимых функций системы.

Толчком к настоящей работе стали беседы с Э. Фишером, в частности его вопрос о минимальной базе лагранжевых областей родов. Относящиеся к этому исследования, которые привели к построению уравнений с заданной группой (ср. цитированное сообщение о ``Полях рациональных функций'', часть 2), оставлены для дальнейшей публикации.

Существенное расширение этой работы по сравнению с первоначальным сообщением состоит в том, что теоремы о базах, сформулированные там только для полей или специальных интегральных областей, здесь перенесены на произвольные системы. Этот перенос удается простым образом благодаря понятию наименьшего поля, содержащего произвольную систему, как обобщению поля частных интегральной области (\S\ 7). К образованию этого понятия меня привело то, что К. Генцельт, узнав рациональную базу полей, смог доказать существование рациональной базы для произвольных систем обходным путем через гильбертову модульную базу. Также к теореме об интегральной базе регулярных систем меня привело замечание К. Генцельта, что рассуждения, примененные мной к интегральным областям, справедливы для произвольных систем, подчиняющихся тем же предположениям; кроме того, я обязана Генцельту еще несколькими отдельными небольшими замечаниями.

Наконец, отметим, что в случае одной неопределенной рациональная база и минимальная база полей имеются у Э. Штайница в его ``Algebraische Theorie der Körper'', \S\ 24 (Crelles Journ. 137); при этом область коэффициентов там берется как поле в самом общем смысле. Вопрос, насколько при этих более общих предположениях сохраняются данные здесь теоремы, ниже не затрагивается; во всяком случае доказательства требуют модификации, поскольку, например, уже средства из теории функциональных детерминант отказывают для полей характеристики $p$.

\end{document}
